\documentclass[11pt,a4paper]{article}
\usepackage[margin=27mm]{geometry}
\usepackage{amsmath,amssymb,amsthm,fontspec,unicode-math,hyperref,booktabs,longtable}
\usepackage{array,calc,fvextra}
\setmainfont{STIXGeneral}
\setmathfont{STIX Two Math}
\setmonofont{Menlo}
\hypersetup{hidelinks}
\setlength{\emergencystretch}{2em}
\providecommand{\tightlist}{\setlength{\itemsep}{0pt}\setlength{\parskip}{0pt}}
\newtheorem{theorem}{Theorem}
\newtheorem{lemma}[theorem]{Lemma}
\newtheorem{corollary}[theorem]{Corollary}
\newcommand{\gen}[1]{\langle #1\rangle}
\newcommand{\Sym}{\operatorname{Sym}}
\newcommand{\Aut}{\operatorname{Aut}}
\newcommand{\Fix}{\operatorname{Fix}}

\title{Stabilizers, extensions and the enumeration of finite geometries}
\author{Human Compression project}
\date{}
\begin{document}
\maketitle
\begin{abstract}
A classification by group actions requires both an equivalence relation and
an admissibility test. We give a common way to extend a partial configuration,
impose a support condition, and count its orbits using its actual stabilizer.
We also prove recognition and exhaustive augmentation for coset geometries
with an arbitrary chamber stabilizer. The applications concern commuting
involutions, Suzuki polyhedra and the residually weakly primitive geometries
of the Hall--Janko group. Their local group structure and exact computations
remain explicit inputs.
\end{abstract}
\section{The information a classification must retain}
Two commuting involutions do not yet specify a regular polyhedron. A third
involution must generate the prescribed group with them and satisfy all
subgroup intersections. Likewise, panel subgroups do not specify a
flag-transitive geometry until incident cosets have common representatives.
The method developed here keeps these tests separate from the orbit reduction.

After fixing a partial configuration, its stabilizer acts on the remaining
choices. Theorem UD.AdmissibleExtensions counts precisely the admissible
choices, including a prescribed field of definition when present. It combines
the elementary extension-orbit correspondence with the classical fixed-point
count and M\"obius inversion. Its role is to expose which hypotheses each
application must prove; no priority is claimed for these classical principles.
The finite-action theorem also records when a permutation image is faithful,
so that an abstract automorphism group is not confused with a normalizer in
one chosen action.

For coset geometries, the common paper contains the complete recognition and
augmentation proofs. The Hall--Janko application supplies the subgroup data,
the actual stabilizer actions and the finite outputs. For Suzuki groups,
normalization and a subfield-position argument identify the support condition
with generation. For symmetric and alternating groups, the transitive
four-group actions supply the multiplicities and parity conditions. These
local arguments are part of their applications rather than assumptions hidden
inside the common theorem.

All groups in the counting statements are finite. The extension-orbit
correspondence and the elementary coset realization are stated without that
restriction. A geometry has a finite type set; incidence is reflexive and
symmetric and two elements of the same type are incident only when equal.
A flag is a pairwise incident set, a chamber has every type, and every flag
must extend to a chamber. Firmness means that a flag missing one type extends
in at least two ways. Residual connectedness concerns the incidence graph of
every residue of rank at least two. The group action is type preserving;
flag transitivity means transitivity on flags of every prescribed type.

A C-group marking is a generating family of distinct nonidentity
involutions $(r_i)_{i\in I}$ such that, with
$G_J=\gen{r_i:i\in J}$ and $G_\varnothing=1$,
\[
 G_J\cap G_K=G_{J\cap K}\qquad(J,K\subseteq I).
\]
For $I=\{0,\ldots,d-1\}$ it is a string marking if
$r_i r_j=r_j r_i$ whenever $|i-j|\ge2$.
The applications retain the standard correspondence between string
C-group markings and regular abstract polytopes. Reversing the ordered
markers corresponds to duality. A rank-three marking is therefore a
regular polyhedron only after generation and all intersections have been
established; endpoint orders alone do not give this conclusion.

\section{Finite actions and
extensions}\label{finite-actions-and-extensions}

\textbf{Theorem UD.Orbits (orbit data and normalizers).} Let a finite
group \(A\) act on a finite set \(\Omega\), and choose one
representative \(H\) of each conjugacy class of subgroups that occurs as
a point stabilizer. Let \(m_H\) be the number of orbits isomorphic to
\(A/H\), with \(m_H=0\) for all omitted classes. Then the numbers
\(m_H\) determine the \(A\)-set up to equivariant isomorphism, and \[
\ker(A\curvearrowright\Omega)=\bigcap_{m_H>0}\operatorname{core}_A(H).
\tag{5.3}
\] Its equivariant permutation group is abstractly \[
C_{\operatorname{Sym}(\Omega)}(A)\cong
\prod_H\bigl((N_A(H)/H)\wr\operatorname{Sym}(m_H)\bigr).
\tag{5.4}
\] If the action \(\rho:A\hookrightarrow\operatorname{Sym}(\Omega)\) is
faithful, conjugation gives an exact sequence \[
1\longrightarrow C_{\operatorname{Sym}(\Omega)}(\rho(A))
\longrightarrow N_{\operatorname{Sym}(\Omega)}(\rho(A))
\longrightarrow
\{\alpha\in\operatorname{Aut}(A):m_{\alpha(H)}=m_H\text{ for all }H\}
\longrightarrow1.
\tag{5.5}
\] No splitting of this sequence is asserted.

\emph{Proof.} Every orbit is equivariantly \(A/H\), where \(H\) is the
stabilizer of a selected point, and two such transitive sets are
isomorphic exactly when the stabilizers are conjugate. This proves the
classification by multiplicities and the kernel formula.

An equivariant permutation preserves the isomorphism type of each orbit.
For one orbit \(A/H\), an equivariant map is determined by the image of
\(H\). That image \(nH\) must be fixed by \(H\), and bijectivity makes
its stabilizer equal to \(H\); hence \(n\in N_A(H)\). Conversely each
such \(n\) gives a permutation \(aH\mapsto anH\). Composition
corresponds to multiplication in the opposite order, so inversion
identifies this permutation group with \(N_A(H)/H\). For \(m_H\) copies,
one may independently choose such a map in every copy and permute the
copies, yielding the wreath product. Different isomorphism types are
independent, proving (5.4).

Faithfulness makes conjugation by a normalizing permutation \(\tau\)
induce a unique automorphism \(\alpha\) of \(A\), with \[
\tau\rho(a)\tau^{-1}=\rho(\alpha(a)).
\] If a point has stabilizer \(H\), its image under \(\tau\) has
stabilizer \(\alpha(H)\); therefore the multiplicities must be
invariant. The kernel of this conjugation map is the centralizer.
Conversely, suppose \(\alpha\) preserves the multiplicities. Pair each
orbit with a target orbit whose stabilizer class is \(\alpha(H)\), and
choose base points \(x,y\) with actual stabilizers \(H,\alpha(H)\),
adjusting within the target orbit if needed. Define
\(\tau(ax)=\alpha(a)y\). The stabilizer equality makes this well-defined
and bijective on that pair of orbits, and the paired definitions combine
to a permutation of \(\Omega\). It satisfies the displayed intertwining
equation, proving surjectivity and exactness. \(\square\)

\textbf{Proposition UD.ExtensionOrbits (classification above a retained
base).} Let a group \(B\) act on sets \(Z,E\), and let \(p:Z\to E\) be
equivariant. For one representative \(e\) of each orbit in
\(B\backslash E\), there is a natural bijection \[
B\backslash Z\ \cong\ \coprod_{e\in B\backslash E}B_e\backslash p^{-1}(e).
\tag{5.6}
\]

\emph{Proof.} Translate any \(z\in Z\) so that \(p(z)\) is the selected
representative of its orbit. If two translated points over that
representative are \(B\)-equivalent, an element taking one to the other
fixes their image \(e\), so belongs to \(B_e\). Conversely a
\(B_e\)-equivalence is already a \(B\)-equivalence. These two
observations prove surjectivity and injectivity of (5.6). \(\square\)

\section{Admissible extensions with prescribed support}
Let a finite group $A$ act on finite sets $Z$ and $E$, and let
$p:Z\to E$ be equivariant. Let $R\subseteq Z$ be invariant; it is the set
passing the application's admissibility tests. Suppose a finite poset $L$
with greatest element $1$ has an $A$-action by order automorphisms and an
equivariant map $s:R\to L$. Write $R_l=\{z\in R:s(z)\le l\}$.
The individual sets $R_l$ need not be invariant. Choose one $e$ in each
$A$-orbit in $p(R)$, and put $R(e)=R\cap p^{-1}(e)$.

\textbf{Theorem UD.AdmissibleExtensions.}
The equivalence classes of admissible extensions of exact support $1$ are
in bijection with the disjoint union of the $A_e$-orbits on
$\{z\in R(e):s(z)=1\}$. Their number is
\[
 \sum_{e\in A\backslash p(R)}\frac1{|A_e|}
 \sum_{a\in A_e}\sum_{l\in L}\mu(l,1)
       |\Fix_{R(e)}(a)\cap R_l|,
\]
where $\mu$ is the M\"obius function of $L$.

\emph{Proof.}
The greatest element is fixed by every order automorphism, so the exact
support condition is invariant. Apply UD.ExtensionOrbits to the restriction
of $p$ to that subset. Its fiber equivalence is exactly the $A_e$-action:
an element carrying one extension over $e$ to another necessarily fixes $e$.
An empty exact-support fiber contributes zero, so representatives may be
taken from all of $p(R)$.

For each $z\in R$, the defining inverse identity of the incidence-algebra
M\"obius function gives
\[
 \sum_{s(z)\le l\le1}\mu(l,1)=
 \begin{cases}1,&s(z)=1,\\0,&s(z)<1.\end{cases}
\]
Sum over the points of $R(e)$ fixed by a particular $a\in A_e$. This gives
the inner sum in the formula. Finally, for a finite group $H$ acting on a
finite set $Y$, count the pairs $(h,y)$ with $hy=y$. Their number is both
$\sum_h|\Fix_Y(h)|$ and $\sum_y|H_y|$. Each orbit contributes $|H|$ to the
second sum by the orbit--stabilizer identity. Division by $|H|$ gives the
number of orbits. Use this identity with $H=A_e$ and the exact-support
fiber, and sum over the representatives $e$. \(\square\)

\textbf{Corollary UD.ExactSupportOrbits.}
If $A$ acts on a finite set $X$ with an equivariant support map to $L$,
the number of exact-top-support orbits is
\[
 \frac1{|A|}\sum_{a\in A}\sum_{l\in L}
       \mu(l,1)|\Fix_X(a)\cap\{x:s(x)\le l\}|.
\]
\emph{Proof.} Take $E$ to be a one-point set and $R=Z=X$.
\(\square\)

Taking a one-element support poset instead recovers counting in ordinary
extension fibers. For an involution $D$ on a finite set $Y$, it gives
$(|Y|+|\Fix_Y(D)|)/2$ orbits. A division by two is therefore justified
without a correction term only after freeness has been proved.

\textbf{Lemma UD.OddOrbits.}
Let a finite group $A$ of odd order act on $X$. Let $D$ be a fixed-point-free
involution of $X$ commuting with this action. Then $D$ has no fixed orbit
in $A\backslash X$.

\emph{Proof.} If $Dx=ax$, commutation and $D^2=1$ give $a^2x=x$.
The order of $a$ is odd, so $a$ is a power of $a^2$ and fixes $x$.
Then $Dx=x$, a contradiction. \(\square\)

For the divisor poset of an integer $N$, support often means the least
field $\mathbb F_{2^d}$ containing all coordinates, where $d\mid N$.
The application must establish that its subfield sets are exactly the
support sublevel sets. If $A(d)$ counts data defined over that field,
the number of data of exact degree $N$ is
\[
 E(N)=\sum_{d\mid N}\mu(N/d)A(d).
\]
If Frobenius acts freely on them with orbit length $N$, division by $N$
counts its orbits. For odd $N$, a commuting fixed-point-free involution
then permits another division by two, by UD.OddOrbits. These are conditional
consequences of the common theorem, not substitutes for the generation and
stabilizer calculations in the applications.

\section{Cosets and flags}
For subgroups \(P_i\leq G\), typed coset Helly means that every pairwise
intersecting family of right cosets \(P_i g_i\), with distinct indices,
has a common element. This property is separate from IP. It supplies a
common representative for an incident flag, whereas IP determines the
stabilizer of a flag already represented in this way.

\textbf{Lemma UD.Helly (anchored criterion).} A finite indexed subgroup
family \((P_i)_{i\in I}\) has typed coset Helly if and only if the
following condition holds. For every \(i\in I\), every
\(J\subseteq I\setminus\{i\}\) with \(|J|\geq2\), and any selected
\(j_0\in J\), each \(g\in P_{j_0}\) satisfying \(g\in P_iP_j\) for all
\(j\in J\) admits \(k\in\bigcap_{j\in J}P_j\) with \(P_i g=P_i k\).

\emph{Proof.} The cosets \(P_i g\) and \(P_j\) intersect precisely when
\(g\in P_iP_j\). Thus Helly gives a common element \(k\), proving
necessity. Conversely, induct on the number of cosets in a pairwise
intersecting family. A family of at most two already has a common
element. For a larger family, the induction hypothesis gives a common
element of all but one coset. Right multiplication by its inverse turns
those cosets into the base cosets \(P_j\). The remaining coset meets
\(P_{j_0}\), so it has a representative \(g\in P_{j_0}\). Its
intersections with the other base cosets are precisely the product
conditions in the statement. The criterion supplies \(k\), and undoing
the right translation gives a common element of the original family.
\(\square\)

\section{Panel systems with a nontrivial chamber
stabilizer}\label{panel-systems-with-a-nontrivial-chamber-stabilizer}

\textbf{Lemma UD.CosetRealization.} Let a group \(G\) act
type-preservingly and flag-transitively on an incidence geometry of
finite rank, and let \((x_i)_{i\in I}\) be a base chamber. Write
\(H_i=\operatorname{Stab}_G(x_i)\) and use a right action. Then \[
H_i g\longmapsto x_i^g
\] is a type-preserving \(G\)-equivariant isomorphism from the subgroup
coset system to the given geometry.

\textbf{Proof.} Transitivity on one-object flags makes the map onto each
type. It is well defined and injective because \(x_i^g=x_i^h\) holds
exactly when \(gh^{-1}\in H_i\), which is equivalent to \(H_i g=H_i h\).
If two cosets share a representative \(k\), their images are
\(x_i^k,x_j^k\) and are incident, as the base objects are incident.
Conversely, if two images are incident, transitivity on flags of their
two types supplies a \(k\) carrying the base pair to them. The injective
object correspondence then writes both cosets with that common
representative. The map therefore preserves and reflects incidence.
\(\square\)

Let \(G\) be finite, let \(I\) be a finite type set, and let each panel
subgroup \(P_i\) contain \(B\). Put \(K_\varnothing=B\) and
\(K_S=\langle P_i:i\in S\rangle\) for nonempty \(S\), with \(K_I=G\).
This marking by subgroups is different from a marking by individual
involutions.

\textbf{Theorem UD.PanelRecognition.} Suppose \[
K_S\cap K_T=K_{S\cap T}\quad(S,T\subseteq I),\qquad [P_i:B]\ge2,
\tag{E4}
\] and suppose the maximal parabolics \(H_i=K_{I\setminus\{i\}}\) have
typed coset Helly. Then their coset incidence system is a firm
flag-transitive residually connected geometry with chamber stabilizer
\(B\). Conversely every finite-rank firm flag-transitive residually
connected coset geometry is obtained this way by taking the panel
stabilizers.

\textbf{Proof.} Equation (E4) identifies the stabilizer of every partial
base flag with the corresponding \(K_S\). UD.Helly supplies a common
group element in the cosets of every flag, and hence flag transitivity
and identification of the residues with the corresponding coset systems
in \(K_S\). Each rank-one residue has \([P_i:B]\) elements, giving
firmness. A rank-two residue is connected because its two panel
subgroups generate \(K_{\{i,j\}}\). For larger rank, all panel subgroups
occur in proper deletions, so the residue group is generated by its
maximal parabolics. The usual coset-path argument proves connectedness:
paths multiply successive elements of those subgroups, and conversely
their products give paths.

Conversely let \(K_S\) stabilize the complementary base flag. Stabilizer
intersections give (E4). In a connected residue, the subgroup generated
by its maximal parabolics has exactly the orbit reachable by coset
paths, so it is the whole residue group. Induction on \(|S|\)
reconstructs \(K_S\) from the panels, starting with \(K_\varnothing=B\).
Flag transitivity gives typed Helly by moving a flag to the base flag.
\(\square\)

The source's remaining axioms are separate, exact local tests. RWPRI
means that for every nonempty \(S\) there is an \(i\in S\) with
\(K_{S\setminus\{i\}}\) maximal proper in \(K_S\). Local
two-transitivity means that \(P_i\) has two orbits on ordered pairs in
\(B\backslash P_i\); equivalently \(B\) has two orbits on that coset
set. The source's geometric \((\mathrm{IP})_2\) means that every
rank-two residue is a partial linear space or a generalized digon. It is
not the subgroup identity (E4).

For the last test let \(P,Q\) be two panels, \(P\cap Q=B\) and
\(K=\langle P,Q\rangle\). The coset graph is a digon exactly when
\(|P||Q|=|K||B|\). Otherwise local two-transitivity makes a single
common-neighbour count sufficient. Choose \(x\in Q\setminus B\).
Neighbours of \(P\) are the \(Q\)-cosets represented by a right
transversal of \(B\) in \(P\); neighbours of \(Px\) are their translates
by \(x\). The stabilizer \(Q\) is two-transitive on its incident
\(P\)-type points, so all distinct points having a common \(Q\)-type
neighbour have the same number of common neighbours. This number must be
one. Pairs with no common neighbour already satisfy partial linearity.
This argument also applies when a local action has a nontrivial kernel.

\subsection{Complete augmentation through maximal
residues}\label{complete-augmentation-through-maximal-residues}

\textbf{Lemma UD.Transport.} Suppose representatives of all ambient
\(G\)-conjugacy classes of subgroups \(K\) and all \(K\)-conjugacy
classes of maximal subgroups \(H_0<K\) are given. For a specified \(H\),
choose any transporter \(c\) with \(H_0^c=H\) when one exists. All
maximal overgroups of \(H\) in that inclusion class are precisely
\((K^c)^n\) with \(n\in N_G(H)\).

\textbf{Proof.} Conjugate any maximal inclusion \(H<L\) to a
representative \(K\) of the class of \(L\). Its image of \(H\) is
\(K\)-conjugate to one of the listed \(H_0\). Comparing the resulting
transporter with \(c\) leaves an element normalizing \(H\). Conversely
each displayed conjugate contains \(H\) maximally. Taking the union over
all inclusion classes gives every maximal overgroup. \(\square\)

Let \(\mathcal P(B)\) be the finite set of overgroups \(P\) of \(B\)
such that \(B\) is maximal proper in \(P\) and \(P\) acts
two-transitively on \(B\backslash P\). UD.Transport computes this entire
set by retaining only the relevant local inclusions. The normalizer
\(N_G(B)\) acts on \(\mathcal P(B)\), and its permutation image has
exactly the same configuration orbits. We retain the actual subgroups
behind the permutation labels, so kernel reduction does not change any
geometric test.

\textbf{Theorem UD.MaximalAugmentation.} Start with all rank-one panels
modulo \(N_G(B)\). For every accepted panel tuple, let \(H\) be its
generated join. For every maximal overgroup \(K\) of \(H\), try every
\(P\in\mathcal P(B)\) with \(P\le K\) and \(P\nsubseteq H\), modulo the
pointwise stabilizer of the old panel tuple. Apply UD.PanelRecognition
and all the local source tests, then remove duplicates under \(N_G(B)\)
and type permutations. Iterating until the next level is empty gives
every target geometry with Borel \(B\) and acting group a subgroup of
\(G\).

\textbf{Proof.} For a target rank-\(r\) geometry, RWPRI supplies at
least one maximal parabolic \(H\) in its acting group \(K\). Deleting
its complementary panel gives a rank-\((r-1)\) residue with the same
Borel and all inherited source axioms. Induction puts its panel tuple in
a previously enumerated orbit. Its remaining panel lies inside \(K\) and
outside \(H\), and UD.Transport includes \(K\) among the maximal
overgroups to be tried. UD.ExtensionOrbits shows that taking
pointwise-stabilizer representatives loses no extension; a smaller
stabilizer than the full setwise one merely permits duplicates. Final
orbit tests remove those duplicates. Conversely every accepted tuple
satisfies exactly the recognition and local conditions.

By (E4) and proper panels, adjoining a new type strictly enlarges the
generated join: if \(P_i\) lay in \(K_S\) for \(i\) outside \(S\), then
\(P_i\cap K_S=B\) would force \(P_i=B\). Hence every tuple gives a chain
of strict subgroup inclusions from \(B\). Its rank is at most
\(\Omega([G:B])\), the number of prime factors with multiplicity. Thus a
run through the next level must reach an empty level. The certificate
requires that emptiness, not merely a chosen numeric cap. Finally take
\(B\) over all \(G\)-conjugacy classes of subgroups. \(\square\)


\section*{Attribution and use}
The orbit classification, centralizer description, fixed-point count and
M\"obius inversion are classical. The common extension correspondence and
the arbitrary-chamber-stabilizer recognition and augmentation proofs are
reused from the previously reviewed project proof bank (U.Orbits,
U.ExtensionOrbits and the enumeration manuscript, Sections 2 and 5).
The combined counting interface removes separate invocations of those
principles; it does not establish new historical classification results.
Coset geometry and the equivalence between string C-groups and abstract
regular polytopes retain their established meanings; the latter correspondence
is an external input in the polyhedron applications.

The full common paper is charged once within this topic. Its overlap with
other topic theories does not reduce that charge. Complete local proofs,
required programs and the separately identified diagnostics accompany the
applications. No numerical compression factor is assigned by this edition.
\end{document}
